% -*- coding: utf-8 -*-
% !TEX program = xelatex
\documentclass[plain,noanswer,medmath]{../../template/jnuexam}
\usepackage{../../template/kysx}

\begin{document}


\examtitle{1992}{5}


\loadproblems{4}{1992P4}%载入试卷四


\makepart{填空题}{本题共5小题，每小题3分，满分15分}


\begin{problem}
设$f(t) = \lim_{x \to \infty} t\left(\frac{x + t}{x - t} \right)^x$，则$f'(t) =$ \fillin{}．
\end{problem}


\useproblem{4}{1}{1}%【同试卷四第一(1)题】


\begin{problem}
设$f(x) = \sin x$，$f[\phi(x)] = 1 - x^2$，则$\phi (x) =$\fillin{}；其定义域为 \fillin{}．
\end{problem}


\begin{problem}
矩阵$\begin{pmatrix}
  1 & 1 & 1 & 1 \\
  1 & 1 & 1 & 1 \\
  1 & 1 & 1 & 1 \\
  1 & 1 & 1 & 1
\end{pmatrix}$的非零特征值是 \fillin{}．
\end{problem}


\begin{problem}
设对于事件$A$,$B$,$C$，有$P(A) = P(B) = P(C) = \frac{1}{4}$，$P(AB) = P(BC) = 0$，
$P(AC) = \frac{1}{8}$，则$A$,$B$,$C$三个事件中至少出现一个的概率为 \fillin{}．
\end{problem}


\makepart{选择题}{本题共5小题，每小题3分，满分15分}


\useproblem{4}{2}{1}%【同试卷四第二(1)题】


%【类似试卷四第二(2)题】
\begin{problem}
当$x \to 0$时，下列四个无穷小量中，哪一个是比其它三个更高阶的无穷小量? \pickout{}
\begin{abcd}
\item $x^2$．     
\item $1 - \cos x$．     
\item $\sqrt{1 - x^2} - 1$．     
\item $x - \sin x$．
\end{abcd}
\end{problem}


\begin{problem}
设$A$, $B$, $A + B$, $A^{-1} + B^{-1}$均为$n$阶可逆矩阵，则$(A^{-1} + B^{-1})^{-1}$等于\pickout{}
\begin{abcd}
\item $A^{-1} + B^{-1}$．
\item $A + B$．
\item $A(A + B)^{-1} B$．
\item $(A + B)^{-1}$．
\end{abcd}
\end{problem}


\begin{problem}
设$\alpha_1$, $\alpha_2$, $\cdots$, $\alpha_m$均为$n$维向量，那么下列结论正确的是\pickout{}
\begin{abcd}
\item 若$k_1\alpha_1 + k_2\alpha_2 + \cdots + k_m \alpha_m = 0$，
      则$\alpha_1$, $\alpha_2$, $\cdots$, $\alpha_m$线性相关．
\item 若对任意一组不全为零的数$k_1$, $k_2$, $\cdots$, $k_m$，
      都有$k_1\alpha_1 + k_2\alpha_2 + \cdots + k_m \alpha_m \ne 0$，
      则$\alpha_1$, $\alpha_2$, $\cdots$, $\alpha_m$线性无关．
\item 若$\alpha_1$, $\alpha_2$, $\cdots$, $\alpha_m$线性相关，
      则对任意一组不全为零的数$k_1$, $k_2$, $\cdots$, $k_m$，
      都有$k_1\alpha_1 + k_2\alpha_2 + \cdots + k_m \alpha_m = 0$．
\item 若$0\alpha_1 + 0\alpha_2 + \cdots + 0\alpha_m = 0$，
      则$\alpha_1$, $\alpha_2$, $\cdots$, $\alpha_m$线性无关．
\end{abcd}
\end{problem}


\useproblem{4}{2}{4}%【同试卷四第二(4)题】


\makepart{}{本题满分5分}%【类似试卷四第三题】

\begin{problem}
求极限$\lim_{x \to 1} \frac{\ln \cos(x - 1)}{1 - \sin \frac{\pi x}{2}}$．
\end{problem}


\makepart{}{本题满分5分}%【同试卷四第四题】

\useproblem{4}{4}{0}


\makepart{}{本题满分6分}%【类似试卷四第六题】

\begin{problem}
求连续函数$f(x)$，使它满足$\int_0^1f(tx)\dt = f(x) + x\sin x$．
\end{problem}


\makepart{}{本题满分5分}%【同试卷四第五题】

\useproblem{4}{5}{0}


\makepart{}{本题满分6分}

\begin{problem}
设生产某产品的固定成本为$10$，而当产量为$x$时的边际成本函数为$MC = - 40 - 20x + 3x^2$，
边际收入函数为$MR = 32 + 10x$，试求：\par
\begin{enumerate*}
\item 总利润函数；
\item 使总利润最大的产量．
\end{enumerate*}
\end{problem}


\makepart{}{本题满分9分}

\begin{problem}
求证：方程$x + p + q\cos x = 0$恰有一个实根，其中$p$,$q$为常数，且$0 < q < 1$．
\end{problem}


\makepart{}{本题满分7分}

\begin{problem}
给定曲线$y = \frac{1}{x^2}$．
\begin{enumerate}
\item 求曲线在横坐标为$x_0$的点处的切线方程；
\item 求曲线的切线被两坐标轴所截线段的最短长度．
\end{enumerate}
\end{problem}


\makepart{}{本题满分5分}%【类似试卷二第三(3)题】

\begin{problem}
设$A$, $B$为$3$阶方阵，$E$为$3$阶单位阵，满足$AB + E = A^2 + B$，又知$A = \begin{pmatrix}
  1 & 0 & 1 \\
  0 & 2 & 0 \\
  1 & 0 & 1
\end{pmatrix}$，求矩阵$B$．
\end{problem}


\makepart{}{本题满分5分}%【类似试卷四第十题】

\begin{problem}
设线性方程组$\begin{cases}
  x_1 + 2x_2 - 2x_3 = 0, \\
  2x_1 - x_2 + \lambda x_3 = 0, \\
  3x_1 + x_2 - x_3 = 0
\end{cases}$的系数矩阵为$A$，三阶矩阵$B \ne 0$，且$AB = 0$．试求求$\lambda$的值．
\end{problem}


\makepart{}{本题满分6分}

\begin{problem}
已知实矩阵$A = (a_{ij})_{3 \times 3}$满足条件：\par
\begin{enumerate*}
\item $A_{ij} = a_{ij}$（$i,j = 1,2,3$），其中$A_{ij}$是$a_{ij}$的代数余子式；
\item $a_{11} \ne 0$．
\end{enumerate*}\par
计算行列式$|A|$．
\end{problem}


\makepart{}{本题满分7分}%【同试卷四第十二题】

\useproblem{4}{12}{0}


\makepart{}{本题满分7分}%【同试卷四第十三题】

\useproblem{4}{13}{0}


\end{document}
